% GENERATED FILE. Edit canonical lessons, then run npm run generate:books.
% canonical-source-hash: fnv1a64:39172b12210c1139
% canonical-lessons: LA-C446-lectulus, LA-C446-stramentum, LA-C446-teges, LA-C446-cervical, LA-C446-lodix, LA-R446-review-439-446

\chapter{Little bed, Straw bedding, Mat, Pillow, Blanket}
\label{ch:la-little-bed-straw-bedding-mat-pillow-blanket}
\hlchaptermodality{446}

\begin{chapteropening}
\textbf{By the end of this chapter:} \emph{I can describe a bed for the night: a little bed, straw bedding, a mat, a pillow and a blanket.}

Complete the last lesson of chapter 446: I can describe a bed for the night: a little bed, straw bedding, a mat, a pillow and a blanket.
\end{chapteropening}

\section[lectulus, lectulī]{lectulus, lectulī — a little bed, a small couch}

\label{lesson:LA-C446-lectulus}



\begin{quote}
Before the new one: say the Latin for a stable-keeper, then the Latin for hospitable.
\end{quote}

\subsection*{You'll want to know: lectulus, lectulī}
\textbf{lectulus, lectulī} — \textquotedblleft{}a little bed, a small couch\textquotedblright{}.

\begin{quote}
In cubiculō parvō lectulus dūrus est.
\end{quote}

\begin{quote}
\emph{In the small bedroom there is a hard little bed.}
\end{quote}

Second declension, masculine: \emph{lectulus, lectulī m.}, a diminutive of \emph{lectus} 'bed'.

\begin{cousinweb}
\emph{lectulus} is a diminutive of \emph{lectus} 'bed, couch' (Lewis \& Short), which is usually traced to PIE \emph{leg\textsuperscript{h}-} 'to lie down'.
\end{cousinweb}

\begin{grammarlens}[title={What you've built}]
One more word for telling, explaining and describing at B1.
\end{grammarlens}

\subsection*{Your turn}
\begin{itemize}
\raggedright
  \item \emph{Say it:} \emph{lectulus, lectulī}
  \item \emph{Say it:} \emph{lectulus, lectulī}, once more
  \item \emph{Say it:} a sentence of your own with \emph{lectulus, lectulī}
  \item \emph{Recall:} say the Latin for a stable-keeper, then the Latin for hospitable, then say \emph{lectulus, lectulī} again
\end{itemize}

\begin{tcolorbox}[breakable,colback=teal!4,colframe=teal!35!black,title={Before you move on}]
What does lectulus, lectulī mean? (\textquotedblleft{}A little bed, a small couch\textquotedblright{}.) Say it once more.
\end{tcolorbox}
\section[strāmentum, strāmentī]{strāmentum, strāmentī — straw bedding, straw; a covering}

\label{lesson:LA-C446-stramentum}



\begin{quote}
Before the new one: say the Latin for hospitable, then the Latin for a little bed.
\end{quote}

\subsection*{You'll want to know: strāmentum, strāmentī}
\textbf{strāmentum, strāmentī} — \textquotedblleft{}straw bedding, straw; a covering\textquotedblright{}.

\begin{quote}
Pauper in strāmentō dormit, nōn in lectō.
\end{quote}

\begin{quote}
\emph{The poor man sleeps on straw, not in a bed.}
\end{quote}

Second declension, neuter: \emph{strāmentum, strāmentī n.} Straw spread out for sleeping on, or for animals.

\begin{cousinweb}
\emph{strāmentum} is built on \emph{sternō} 'to spread' (Lewis \& Short), from its participle \emph{strātus}, with the suffix \emph{-mentum}: something spread out to lie on.
\end{cousinweb}

\begin{grammarlens}[title={What you've built}]
One more word for telling, explaining and describing at B1.
\end{grammarlens}

\subsection*{Your turn}
\begin{itemize}
\raggedright
  \item \emph{Say it:} \emph{strāmentum, strāmentī}
  \item \emph{Say it:} \emph{strāmentum, strāmentī}, once more
  \item \emph{Say it:} a sentence of your own with \emph{strāmentum, strāmentī}
  \item \emph{Recall:} say the Latin for hospitable, then the Latin for a little bed, then say \emph{strāmentum, strāmentī} again
\end{itemize}

\begin{tcolorbox}[breakable,colback=teal!4,colframe=teal!35!black,title={Before you move on}]
What does strāmentum, strāmentī mean? (\textquotedblleft{}Straw bedding, straw; a covering\textquotedblright{}.) Say it once more.
\end{tcolorbox}
\section[teges, tegetis]{teges, tegetis — a mat, a rush mat; a covering}

\label{lesson:LA-C446-teges}



\begin{quote}
Before the new one: say the Latin for a little bed, then the Latin for straw bedding.
\end{quote}

\subsection*{You'll want to know: teges, tegetis}
\textbf{teges, tegetis} — \textquotedblleft{}a mat, a rush mat; a covering\textquotedblright{}.

\begin{quote}
In caupōnā parvā tantum teges in pavīmentō est.
\end{quote}

\begin{quote}
\emph{In a small inn there is only a mat on the floor.}
\end{quote}

Third declension, feminine: \emph{teges, tegetis f.} A coarse mat of rushes, the poor man's bed in Roman satire.

\begin{cousinweb}
\emph{teges} is built on \emph{tegō} 'to cover' (Lewis \& Short): a covering, a mat.
\end{cousinweb}

\begin{grammarlens}[title={What you've built}]
One more word for telling, explaining and describing at B1.
\end{grammarlens}

\subsection*{Your turn}
\begin{itemize}
\raggedright
  \item \emph{Say it:} \emph{teges, tegetis}
  \item \emph{Say it:} \emph{teges, tegetis}, once more
  \item \emph{Say it:} a sentence of your own with \emph{teges, tegetis}
  \item \emph{Recall:} say the Latin for a little bed, then the Latin for straw bedding, then say \emph{teges, tegetis} again
\end{itemize}

\begin{tcolorbox}[breakable,colback=teal!4,colframe=teal!35!black,title={Before you move on}]
What does teges, tegetis mean? (\textquotedblleft{}A mat, a rush mat; a covering\textquotedblright{}.) Say it once more.
\end{tcolorbox}
\section[cervīcal, cervīcālis]{cervīcal, cervīcālis — a pillow, a bolster}

\label{lesson:LA-C446-cervical}



\begin{quote}
Before the new one: say the Latin for straw bedding, then the Latin for a mat.
\end{quote}

\subsection*{You'll want to know: cervīcal, cervīcālis}
\textbf{cervīcal, cervīcālis} — \textquotedblleft{}a pillow, a bolster\textquotedblright{}.

\begin{quote}
Hospes cervīcal sub caput pōnit et dormit.
\end{quote}

\begin{quote}
\emph{The guest puts a pillow under his head and sleeps.}
\end{quote}

Third declension, neuter: \emph{cervīcal, cervīcālis n.} Literally a thing for the neck.

\begin{cousinweb}
\emph{cervīcal} is built on \emph{cervīx, cervīcis} 'neck' (EtymDB; Lewis \& Short): a support for the neck. English \emph{cervical} comes from \emph{cervīx}.
\end{cousinweb}

\begin{grammarlens}[title={What you've built}]
One more word for telling, explaining and describing at B1.
\end{grammarlens}

\subsection*{Your turn}
\begin{itemize}
\raggedright
  \item \emph{Say it:} \emph{cervīcal, cervīcālis}
  \item \emph{Say it:} \emph{cervīcal, cervīcālis}, once more
  \item \emph{Say it:} a sentence of your own with \emph{cervīcal, cervīcālis}
  \item \emph{Recall:} say the Latin for straw bedding, then the Latin for a mat, then say \emph{cervīcal, cervīcālis} again
\end{itemize}

\begin{tcolorbox}[breakable,colback=teal!4,colframe=teal!35!black,title={Before you move on}]
What does cervīcal, cervīcālis mean? (\textquotedblleft{}A pillow, a bolster\textquotedblright{}.) Say it once more.
\end{tcolorbox}
\section[lōdīx, lōdīcis]{lōdīx, lōdīcis — a blanket, a coverlet}

\label{lesson:LA-C446-lodix}



\begin{quote}
Before the new one: say the Latin for a mat, then the Latin for a pillow.
\end{quote}

\subsection*{You'll want to know: lōdīx, lōdīcis}
\textbf{lōdīx, lōdīcis} — \textquotedblleft{}a blanket, a coverlet\textquotedblright{}.

\begin{quote}
Nox frīgida est, sed lōdīx tenuis est.
\end{quote}

\begin{quote}
\emph{The night is cold, but the blanket is thin.}
\end{quote}

Third declension, feminine: \emph{lōdīx, lōdīcis f.} A wool blanket for a bed.

\begin{cousinweb}
The origin of \emph{lōdīx} is not settled: Lewis \& Short give the sense, a blanket or coverlet, but no source, and EtymDB does not list it.
\end{cousinweb}

\begin{grammarlens}[title={What you've built}]
One more word for telling, explaining and describing at B1.
\end{grammarlens}

\subsection*{Your turn}
\begin{itemize}
\raggedright
  \item \emph{Say it:} \emph{lōdīx, lōdīcis}
  \item \emph{Say it:} \emph{lōdīx, lōdīcis}, once more
  \item \emph{Say it:} a sentence of your own with \emph{lōdīx, lōdīcis}
  \item \emph{Recall:} say the Latin for a mat, then the Latin for a pillow, then say \emph{lōdīx, lōdīcis} again
\end{itemize}

\begin{tcolorbox}[breakable,colback=teal!4,colframe=teal!35!black,title={Before you move on}]
What does lōdīx, lōdīcis mean? (\textquotedblleft{}A blanket, a coverlet\textquotedblright{}.) Say it once more.
\end{tcolorbox}
\section[(dialogue)]{Words from chapters 439–446 — a review}

\label{lesson:LA-R446-review-439-446}



\begin{quote}
Say the Latin for to put in and to step out.
\end{quote}

\subsection*{Your turn}
Say these aloud:

\begin{itemize}
\raggedright
  \item to put in, to step out, to haul ashore, to launch, navigable
  \item to toss on the waves, to be seasick, to carry down, to take across, to swim to safety
  \item a harbour, a dockyard, a tidal inlet, a lighthouse, a headland
  \item the open sea, the high sea, the level sea, of the sea, coastal
  \item a sandbank, a billow, a shallow, a strait, to sail past
  \item a dead calm, calm weather, a peninsula, of the shore, to sail out
  \item a roadside inn, to lodge, to spend the night, a stable-keeper, hospitable
  \item a little bed, straw bedding, a mat, a pillow, a blanket
\end{itemize}

\begin{tcolorbox}[breakable,colback=teal!4,colframe=teal!35!black,title={Before you move on}]
To put in? (\textbf{appellō, appellere}.) To step out? (\textbf{ēgredior, ēgredī}.) To haul ashore? (\textbf{subdūcō, subdūcere}.) To launch? (\textbf{dēdūcō, dēdūcere}.) Navigable? (\textbf{nāvigābilis, nāvigābile}.) To toss on the waves? (\textbf{flūctuō, flūctuāre}.) To be seasick? (\textbf{nauseō, nauseāre}.) To carry down? (\textbf{dēferō, dēferre}.) To take across? (\textbf{trāiciō, trāicere}.) To swim to safety? (\textbf{ēnatō, ēnatāre}.) A harbour? (\textbf{portus, portūs}.) A dockyard? (\textbf{nāvāle, nāvālis}.) A tidal inlet? (\textbf{aestuārium, aestuāriī}.) A lighthouse? (\textbf{pharus, pharī}.) A headland? (\textbf{prōmunturium, prōmunturiī}.) The open sea? (\textbf{salum, salī}.) The high sea? (\textbf{pelagus, pelagī}.) The level sea? (\textbf{aequor, aequoris}.) Of the sea? (\textbf{marīnus, marīna, marīnum}.) Coastal? (\textbf{maritimus, maritima, maritimum}.) A sandbank? (\textbf{syrtis, syrtis}.) A billow? (\textbf{flūctus, flūctūs}.) A shallow? (\textbf{vadum, vadī}.) A strait? (\textbf{fretum, fretī}.) To sail past? (\textbf{praetervehor, praetervehī}.) A dead calm? (\textbf{malacia, malaciae}.) Calm weather? (\textbf{tranquillitās, tranquillitātis}.) A peninsula? (\textbf{paenīnsula, paenīnsulae}.) Of the shore? (\textbf{lītoreus, lītorea, lītoreum}.) To sail out? (\textbf{ēnāvigō, ēnāvigāre}.) A roadside inn? (\textbf{dēverticulum, dēverticulī}.) To lodge? (\textbf{dēversor, dēversārī}.) To spend the night? (\textbf{pernoctō, pernoctāre}.) A stable-keeper? (\textbf{stabulārius, stabulāriī}.) Hospitable? (\textbf{hospitālis, hospitāle}.) A little bed? (\textbf{lectulus, lectulī}.) Straw bedding? (\textbf{strāmentum, strāmentī}.) A mat? (\textbf{teges, tegetis}.) A pillow? (\textbf{cervīcal, cervīcālis}.) A blanket? (\textbf{lōdīx, lōdīcis}.)
\end{tcolorbox}
